\chapter{Logic and Proof}

\begin{center}\itshape
If controversies were to arise, there would be no more need of disputation between two philosophers than between\\
two accountants. For it would suffice to take their pencils in their hands [\dots] and to say to each other [\dots]: Let us calculate.
\upshape(Leibniz, about 1685)
\end{center}

In 1847 a schoolmaster from Lincoln, the son of a shoemaker, who had taught himself mathematics from library books,
published a small pamphlet with an outrageous claim: thinking obeys algebra. In \emph{The Laws of Thought} (1854) George Boole writes $x$ for ``white
things'' and $y$ for ``the sheep'', so that $xy$ means ``white sheep''; and then $xx=x$, since the white white things are
just the white things. That equation holds for no number except $0$ and $1$. So logic, he concludes, is the algebra of the two numbers $0$ and $1$. Leibniz had
dreamed of such a calculus a century and a half earlier, and the philosophers of Boole's day mostly ignored it.

Ninety years later, a student at MIT named Claude Shannon noticed that the relays of a telephone exchange obey
Boole's laws: two switches in series are an \emph{and}, two in parallel an \emph{or}. His master's thesis of 1937,
perhaps the most consequential master's thesis ever written, turned Boole's algebra into the design method of
every digital circuit since. Today the same algebra runs the other way too: we ask computers to \emph{reason}.
In 2016 Marijn Heule, Oliver Kullmann and Victor Marek asked whether the integers $1,2,\dots,n$ can be coloured red and
blue so that no Pythagorean triple $a^2+b^2=c^2$ is all one colour. Nobody knew. A satisfiability solver, running on a
supercomputer for two days, found that it can be done up to $n=7824$ and is impossible for $n=7825$, and it produced a
proof of impossibility two hundred terabytes long, the longest proof ever written at the time (Heule beat his own
record a year later with two petabytes, for a problem of Chapter~8). Every step of that proof is one application of a
single rule of this chapter.

This chapter tells that story in order. We first show that every truth table is a formula, and that some truth tables
need enormous ones. We give Boolean algebra its abstract form and prove that every finite Boolean algebra is an
algebra of subsets, which is the normal form theorem in disguise. Then we turn to proofs: \emph{resolution}, the single
rule behind the Pythagorean proof, which refutes every unsatisfiable set of clauses; the \emph{compactness theorem},
which lets infinite problems be solved by finite ones; and two fragments where satisfiability is easy, Horn clauses
and clauses of length two. The chapter ends with the language of first-order logic, in which the rest of
mathematics is written.

\begin{think}
Can the following formula be made true, and if so how?
\[
\begin{aligned}{}&(x_1\lor x_2)\land(\neg x_1\lor x_3)\land(\neg x_2\lor\neg x_3)\eqbr\land(\neg x_3\lor x_4)\land(\neg x_4\lor\neg x_1)\land(x_2\lor x_4).\end{aligned}
\]
What if we add the clause $(\neg x_2\lor x_1)$? With $4$ variables you can try all $16$ cases; with $300$ you
could not. Section~\ref{sec:twosat} gives a method that decides such formulas, with any number of variables, in
a time proportional to their length.
\end{think}

\section{Propositional formulas}

Fix a set of \emph{variables} $x_1,x_2,\dots$. \emph{Formulas} are built from the variables and the constants $\top$
(true) and $\bot$ (false) by the rules: if $\varphi$ and $\psi$ are formulas, so are $\neg\varphi$, $(\varphi\land\psi)$,
$(\varphi\lor\psi)$, $(\varphi\to\psi)$ and $(\varphi\leftrightarrow\psi)$. It is best to think of a formula as a
tree, its \emph{parse tree}: the leaves carry variables or constants and the inner nodes carry connectives. The
parentheses only serve to write the tree on a line, and we drop those that are clear from context. Since a
formula is a finite tree, we can define functions on formulas and prove properties of them \emph{by structural
induction}: handle the leaves, then a node given its subtrees.

An \emph{assignment} is a map $\alpha$ from the variables to $\{0,1\}$ ($0$ for false, $1$ for true). It extends to all
formulas by induction: $\alpha(\top)=1$, $\alpha(\bot)=0$, $\alpha(\neg\varphi)=1-\alpha(\varphi)$, $\alpha(\varphi\land\psi)=\min(\alpha(\varphi),\alpha(\psi))$,
$\alpha(\varphi\lor\psi)=\max(\alpha(\varphi),\alpha(\psi))$, $\alpha(\varphi\to\psi)=\max(1-\alpha(\varphi),\alpha(\psi))$, and $\alpha(\varphi\leftrightarrow\psi)=1$ if and only
if $\alpha(\varphi)=\alpha(\psi)$. The value $\alpha(\varphi)$ only depends on the values of $\alpha$ at the variables that occur in $\varphi$ (an
induction on $\varphi$).

\begin{definition}\label{def:sat}
A formula $\varphi$ is \emph{satisfiable} if $\alpha(\varphi)=1$ for some assignment $\alpha$, which is then a \emph{model} of
$\varphi$, and a \emph{tautology} if $\alpha(\varphi)=1$ for all $\alpha$. A set $\Gamma$ of formulas is satisfiable if one
assignment satisfies all its members. We write $\Gamma\models\varphi$, and say that $\varphi$ is a \emph{consequence} of
$\Gamma$, if every model of $\Gamma$ is a model of $\varphi$. Two formulas are \emph{equivalent}, $\varphi\equiv\psi$, if
$\alpha(\varphi)=\alpha(\psi)$ for every $\alpha$.
\end{definition}

The basic equivalences are checked on the four rows of a truth table: commutativity and associativity of $\land$ and
$\lor$, the two distributive laws $\varphi\land(\psi\lor\chi)\equiv(\varphi\land\psi)\lor(\varphi\land\chi)$ and
$\varphi\lor(\psi\land\chi)\equiv(\varphi\lor\psi)\land(\varphi\lor\chi)$, De Morgan's laws $\neg(\varphi\land\psi)\equiv\neg\varphi\lor\neg\psi$ and
$\neg(\varphi\lor\psi)\equiv\neg\varphi\land\neg\psi$, double negation $\neg\neg\varphi\equiv\varphi$, and $\varphi\to\psi\equiv\neg\varphi\lor\psi$. Replacing a
subformula by an equivalent one gives an equivalent formula (induction on the parse tree). Satisfiability and
consequence are linked by a remark that we will use all the time.

\begin{proposition}\label{prop:refute}
$\Gamma\models\varphi$ if and only if $\Gamma\cup\{\neg\varphi\}$ is unsatisfiable. In particular $\varphi$ is a tautology if and only
if $\neg\varphi$ is unsatisfiable.
\end{proposition}
\begin{proof}
A model of $\Gamma\cup\{\neg\varphi\}$ is exactly a model of $\Gamma$ in which $\varphi$ is false. There is none if and only if every
model of $\Gamma$ satisfies $\varphi$.
\end{proof}

So to \emph{prove} a statement it is enough to show that its negation, added to the hypotheses, cannot be
satisfied. This is how SAT solvers are used to verify hardware and software: the property to check is negated, and
the solver either finds a satisfying assignment, which is a bug, or shows there is none, which is a proof.

\section{Normal forms}

A \emph{Boolean function} of $n$ variables is a map $f\colon\{0,1\}^n\to\{0,1\}$; there are $2^{2^n}$ of them. Each formula in
$x_1,\dots,x_n$ \emph{computes} one, $f(a)=\alpha_a(\varphi)$, where $\alpha_a$ sends $x_i$ to $a_i$. A \emph{literal} is a variable or its
negation; for a literal $\ell$ we write $\bar\ell$ for its opposite ($\bar x=\neg x$ and $\overline{\neg x}=x$). A \emph{clause} is a
disjunction of literals, and a formula is in \emph{conjunctive normal form} (CNF) if it is a conjunction of clauses,
in \emph{disjunctive normal form} (DNF) if it is a disjunction of conjunctions of literals. The empty disjunction is
$\bot$ and the empty conjunction is $\top$.

\begin{theorem}\label{thm:normalforms}
Every Boolean function is computed by a formula in DNF and by a formula in CNF.
\end{theorem}
\begin{proof}
For $a\in\{0,1\}^n$, let $m_a=\ell_1\land\dots\land\ell_n$ with $\ell_i=x_i$ if $a_i=1$ and $\ell_i=\neg x_i$ if $a_i=0$. Then $m_a$ is true exactly at
$a$: it is the \emph{minterm} of $a$. The DNF $\bigvee_{f(a)=1}m_a$ is true exactly where $f$ is $1$, so it computes $f$. Applying
this to $1-f$ gives a DNF $\psi$ for $1-f$; then $\neg\psi$ computes $f$, and De Morgan's laws push the negation inside: $\neg\bigvee_a m_a\equiv\bigwedge_a\neg m_a$
and $\neg m_a\equiv\bar\ell_1\lor\dots\lor\bar\ell_n$, a clause. So $f$ is computed by the CNF $\bigwedge_{f(a)=0}(\bar\ell_1\lor\dots\lor\bar\ell_n)$: one clause forbidding each
point where $f$ vanishes.
\end{proof}

A set of connectives is \emph{complete} if every Boolean function is computed by a formula that uses only them (and
variables). Theorem~\ref{thm:normalforms} says that $\{\neg,\land,\lor\}$ is complete; since $\varphi\lor\psi\equiv\neg(\neg\varphi\land\neg\psi)$, so is $\{\neg,\land\}$. The
\emph{Sheffer stroke} or NAND, $\varphi\uparrow\psi=\neg(\varphi\land\psi)$, is complete on its own, because $\neg\varphi\equiv\varphi\uparrow\varphi$ and $\varphi\land\psi\equiv(\varphi\uparrow\psi)\uparrow(\varphi\uparrow\psi)$;
this is why a chip can be built from a single kind of gate. Not every set is complete:

\begin{proposition}\label{prop:monotone}
Call $f$ \emph{monotone} if $a\le b$ coordinatewise implies $f(a)\le f(b)$. Every formula built from variables with $\land$ and
$\lor$ only computes a monotone function. In particular $\{\land,\lor\}$ is not complete.
\end{proposition}
\begin{proof}
By structural induction: a variable is monotone, and if $g$ and $h$ are monotone, so are $\min(g,h)$ and $\max(g,h)$, as
$\min$ and $\max$ are non-decreasing in each argument. The function $\neg x_1$ is not monotone.
\end{proof}

Exercise~\ref{ex:monotone} shows the converse: every monotone function has a formula without negation.

The normal forms of Theorem~\ref{thm:normalforms} can be huge, and sometimes this cannot be helped. The \emph{parity
function} $x_1\oplus\dots\oplus x_n$ is $1$ when an odd number of the $x_i$ are $1$.

\begin{proposition}\label{prop:parity}
Every CNF that computes the parity of $n$ variables has at least $2^{n-1}$ clauses.
\end{proposition}
\begin{proof}
Remove from the CNF any clause containing both a variable and its negation: such a clause is always true. Let $C$ be one
of the remaining clauses; it is false at some assignment $a$ (make each of its literals false, which is possible as no
two of them are opposite). Suppose a variable $x_i$ does not occur in $C$, and let $a'$ be $a$ with the $i$-th bit flipped.
Then $C$ is still false at $a'$, so the CNF is false at $a$ and at $a'$, but the parity of $a$ and $a'$ differ: one of them is
$1$. This is impossible, so every clause contains all $n$ variables, and is then false at exactly one assignment. The
CNF must be false at each of the $2^{n-1}$ assignments of even parity, and each clause accounts for only one of
them.
\end{proof}

Even worse, Exercise~\ref{ex:shannoncount} shows by counting that most Boolean functions of $n$ variables need formulas
of length about $2^n/\log_2n$, whatever the form. But a satisfiability solver does not need an \emph{equivalent} CNF; it only
needs one that is satisfiable exactly when the original formula is. That can always be had cheaply, at the price of
new variables.

\begin{theorem}[Tseitin, 1968]\label{thm:tseitin}
For every formula $\varphi$ with $s$ connectives there is a CNF $\psi$ with at most $4s+1$ clauses, each of at most three
literals, in the variables of $\varphi$ and $s$ new variables, such that every model of $\varphi$ extends in exactly one way to a
model of $\psi$, and every model of $\psi$ restricts to a model of $\varphi$. In particular $\varphi$ and $\psi$ are \emph{equisatisfiable}.
\end{theorem}
\begin{proof}
We may assume that $\varphi$ contains no constant, since $\top$ and $\bot$ can be removed by the equivalences $\chi\land\top\equiv\chi$,
$\chi\land\bot\equiv\bot$, $\neg\top\equiv\bot$ and their analogues for the other connectives, unless $\varphi$ reduces to a constant, a trivial case.
Give each inner node $u$ of the parse tree a new variable $y_u$, meant to hold the value of the subformula at $u$, and for a
leaf let $y_u$ be its variable. For each inner node add clauses saying that $y_u$ is computed correctly from its children
$v,w$:
\begin{itemize}
\item for $u=\neg v$: $(y_u\lor y_v)\land(\neg y_u\lor\neg y_v)$, which says $y_u\leftrightarrow\neg y_v$;
\item for $u=v\land w$: $(\neg y_u\lor y_v)\land(\neg y_u\lor y_w)\land(y_u\lor\neg y_v\lor\neg y_w)$, which says $y_u\leftrightarrow y_v\land y_w$;
\item for $u=v\lor w$: $(y_u\lor\neg y_v)\land(y_u\lor\neg y_w)\land(\neg y_u\lor y_v\lor y_w)$, which says $y_u\leftrightarrow y_v\lor y_w$;
\item for $u=v\to w$: $(y_u\lor y_v)\land(y_u\lor\neg y_w)\land(\neg y_u\lor\neg y_v\lor y_w)$;
\item for $u=v\leftrightarrow w$: the four clauses $(\neg y_u\lor\neg y_v\lor y_w)$, $(\neg y_u\lor y_v\lor\neg y_w)$, $(y_u\lor y_v\lor y_w)$ and
$(y_u\lor\neg y_v\lor\neg y_w)$, which exclude exactly the four assignments of $(y_u,y_v,y_w)$ where $y_u$ is not $y_v\leftrightarrow y_w$.
\end{itemize}
(Each line is the CNF of Theorem~\ref{thm:normalforms} for the function ``$y_u$ equals the value of the connective'', simplified.)
Finally add the one-literal clause $(y_r)$ for the root $r$. That is at most $4s+1$ clauses. Given a model $\alpha$ of $\varphi$, the
clauses for the nodes force, from the leaves up, $y_u=\alpha(\text{subformula at }u)$, which is a model of all clauses including
$y_r$, and it is the only extension. Conversely in a model of $\psi$, the same induction shows that $y_u$ equals the value of
the subformula at $u$, so $y_r=1$ means that $\varphi$ is true.
\end{proof}

So satisfiability of arbitrary formulas reduces, in linear time, to satisfiability of CNFs with clauses of length
at most $3$: the problem \textsc{3-SAT}. Tome~C3 shows that every problem whose solutions can be checked quickly reduces
to it in the same way (the Cook--Levin theorem), and that is why \textsc{3-SAT} is the most studied problem of computer
science.

\section{Boolean algebras}\label{sec:boolean}

Formulas up to equivalence, subsets of a set, and the divisors of $30$ under gcd and lcm all obey the same laws. Here are
the laws.

\begin{definition}\label{def:boolalg}
A \emph{Boolean algebra} is a set $B$ with two binary operations $\land,\lor$, a unary operation $x\mapsto\neg x$ and two elements
$0,1$, such that for all $x,y,z$: $\land$ and $\lor$ are commutative and associative; each distributes over the other;
$x\land1=x$ and $x\lor0=x$; and $x\land\neg x=0$ and $x\lor\neg x=1$.
\end{definition}

Examples: the subsets of a set $X$ with $\cap,\cup$, complement, $\emptyset$ and $X$; the two-element algebra $\{0,1\}$; the
\emph{Lindenbaum algebra} of formulas in $x_1,\dots,x_n$ up to equivalence; the divisors of a squarefree integer $N$ with $\gcd$,
$\operatorname{lcm}$, $d\mapsto N/d$, $1$ and $N$. The axioms come in dual pairs (swap $\land$ with $\lor$ and $0$ with $1$), so every consequence of
the axioms has a dual, also true.

\begin{lemma}\label{lem:boolbasic}
In a Boolean algebra: (a) $x\land x=x$, $x\land0=0$ and $x\land(x\lor y)=x$, and dually; (b) if $x\land y=0$ and $x\lor y=1$ then $y=\neg x$;
in particular $\neg\neg x=x$; (c) De Morgan: $\neg(x\land y)=\neg x\lor\neg y$; (d) the relation $x\le y\iff x\land y=x$ is a partial order, and
$x\land y=x$ if and only if $x\lor y=y$.
\end{lemma}
\begin{proof}
(a) $x=x\land1=x\land(x\lor\neg x)=(x\land x)\lor(x\land\neg x)=(x\land x)\lor0=x\land x$. Then
$x\land0=x\land(x\land\neg x)=(x\land x)\land\neg x=x\land\neg x=0$. And $x\land(x\lor y)=(x\lor0)\land(x\lor y)=x\lor(0\land y)=x\lor0=x$, by the second distributive law
and $0\land y=0$.

(b) $y=y\land1=y\land(x\lor\neg x)=(y\land x)\lor(y\land\neg x)=0\lor(y\land\neg x)=y\land\neg x$, and symmetrically, using $x\lor y=1$ and
$x\land\neg x=0$: $\neg x=\neg x\land(x\lor y)=(\neg x\land x)\lor(\neg x\land y)=\neg x\land y$. So $y=\neg x$. Since $\neg x\land x=0$ and $\neg x\lor x=1$, the same
argument with $\neg x$ in the role of $x$ gives $x=\neg\neg x$.

(c) By (b) it suffices to check that $(x\land y)\land(\neg x\lor\neg y)=0$ and $(x\land y)\lor(\neg x\lor\neg y)=1$. The first is
$(x\land y\land\neg x)\lor(x\land y\land\neg y)=0\lor0$ by distributivity and (a). The second is, by distributivity of $\lor$ over $\land$,
$(x\lor\neg x\lor\neg y)\land(y\lor\neg x\lor\neg y)=1\land1$, using the dual of $x\land0=0$, namely $x\lor1=1$.

(d) If $x\land y=x$ then $x\lor y=(x\land y)\lor y=y$ by the dual absorption law, and conversely if $x\lor y=y$ then
$x\land y=x\land(x\lor y)=x$. Reflexivity is $x\land x=x$; if $x\le y\le x$ then $x=x\land y=y\land x=y$; if $x\le y\le z$ then
$x\land z=(x\land y)\land z=x\land(y\land z)=x\land y=x$.
\end{proof}

For this order, $x\land y$ is the greatest lower bound of $x$ and $y$ and $x\lor y$ their least upper bound (Exercise~\ref{ex:lattice}), $0$ is
the least element and $1$ the greatest. An \emph{atom} is a minimal non-zero element: $a\ne0$, and $0\le x\le a$ implies $x=0$
or $x=a$. In the algebra of subsets, the atoms are the singletons. In the Lindenbaum algebra, they are the classes of
the $2^n$ minterms.

\begin{theorem}[finite Stone representation]\label{thm:stone}
Let $B$ be a finite Boolean algebra and $A$ its set of atoms. The map $\Phi(x)=\{a\in A:a\le x\}$ is an isomorphism from $B$ to
the algebra of subsets of $A$. In particular $|B|=2^{|A|}$.
\end{theorem}
\begin{proof}
Three facts about an atom $a$. (i) For every $x$, $a\land x$ is $0$ or $a$, because $a\land x\le a$. (ii) $a\le x\lor y$ if and only if
$a\le x$ or $a\le y$: indeed $a\land(x\lor y)=(a\land x)\lor(a\land y)$, and by (i) this is $a$ exactly when one of the two terms is $a$.
(iii) $a\le\neg x$ if and only if $a\not\le x$: if $a\land x=0$ then $a=a\land(x\lor\neg x)=a\land\neg x$, so $a\le\neg x$; and if $a\le x$ and $a\le\neg x$
then $a\le x\land\neg x=0$, which is false.

Also $a\le x\land y$ if and only if $a\le x$ and $a\le y$: if $a\land x=a=a\land y$ then $a\land(x\land y)=(a\land x)\land y=a$, and if
$a\land x\land y=a$ then $a\land x=(a\land x\land y)\land x=a\land x\land y=a$, and likewise for $y$. By these facts and (ii), (iii), $\Phi(x\lor y)=\Phi(x)\cup\Phi(y)$, $\Phi(\neg x)=A\setminus\Phi(x)$ and $\Phi(x\land y)=\Phi(x)\cap\Phi(y)$; also
$\Phi(0)=\emptyset$ and $\Phi(1)=A$. So $\Phi$ is a homomorphism. Every non-zero $z$ lies above an atom, since $B$ is finite: take an
element minimal among the non-zero elements below $z$. \emph{Injectivity}: let $y=\bigvee\Phi(x)$, the join of the atoms below $x$ (with
$y=0$ if there are none). Then $y\le x$. If $y\ne x$, then $z=x\land\neg y\ne0$, for otherwise $x=x\land(y\lor\neg y)=(x\land y)\lor0=y$. An atom $a\le z$
satisfies $a\le x$, hence $a\le y$, and $a\le\neg y$, which contradicts (iii). So $x=\bigvee\Phi(x)$ is determined by $\Phi(x)$.
\emph{Surjectivity}: for $S\subseteq A$, let $x=\bigvee S$. By (ii) applied repeatedly, an atom $b$ lies below $x$ if and only if it lies
below some $a\in S$, that is, $b=a$. So $\Phi(x)=S$.
\end{proof}

\begin{beauty}
The Lindenbaum algebra of $n$ variables has $2^n$ atoms, the minterms, so it has $2^{2^n}$ elements: one for each Boolean
function. Stone's theorem is then exactly the disjunctive normal form: an element is the join of the atoms below it,
and a formula is the disjunction of the minterms of the points where it is true. The same theorem, with
topology added, describes infinite Boolean algebras too (Stone, 1936), and it was one of the first great
dualities of modern mathematics.
\end{beauty}

\section{Resolution}\label{sec:resolution}

In 1960 Martin Davis and Hilary Putnam wanted a computer to prove theorems, and they found that the heart of the
matter was propositional. Five years later John Alan Robinson distilled their method into one rule, so simple that it
looks like it cannot possibly be enough. It is enough.

From now on a CNF is viewed as a finite set $S$ of clauses, and a clause as a finite set of literals; the empty clause,
written $\Box$, is false under every assignment, and a set containing it is unsatisfiable. A clause containing a literal
and its opposite is a \emph{tautology}.

\begin{definition}\label{def:resolution}
If $C\cup\{x\}$ and $D\cup\{\neg x\}$ are clauses, the clause $C\cup D$ is their \emph{resolvent} on $x$. A \emph{resolution derivation} from $S$
is a sequence of clauses, each of which is in $S$ or is a resolvent of two earlier ones. A \emph{refutation} of $S$ is a
derivation that ends with $\Box$.
\end{definition}

For instance, from
\[
S=\bigl\{\{x,y\},\{x,\neg y\},\{\neg x,y\},\{\neg x,\neg y\}\bigr\}
\]
(the four ways of being wrong about two bits) we derive $\{x\}$ from the
first two clauses, $\{\neg x\}$ from the last two, and $\Box$ from these. Figure~\ref{fig:resolution} draws a refutation as a tree.

\begin{figure}[ht]
\centering
\fig{ch01-resolution}
\caption{A resolution refutation of the second formula of the opening question, the one with seven clauses. Each inner
node is the resolvent of its two parents on the variable written beside it; the root is the empty clause.}
\label{fig:resolution}
\end{figure}

\begin{lemma}[soundness]\label{lem:sound}
A resolvent is a consequence of the two clauses it comes from. Hence if $S$ has a refutation, $S$ is unsatisfiable.
\end{lemma}
\begin{proof}
Let $\alpha$ satisfy $C\cup\{x\}$ and $D\cup\{\neg x\}$. If $\alpha(x)=1$, then $\neg x$ is false, so some literal of $D$ is true; if $\alpha(x)=0$, some
literal of $C$ is true. Either way $\alpha$ satisfies $C\cup D$. By induction along a derivation, every model of $S$ satisfies every
derived clause, and no assignment satisfies $\Box$.
\end{proof}

The converse is the remarkable part: a single rule is enough to refute every unsatisfiable set of clauses. The proof
eliminates the variables one at a time, an idea of Davis and Putnam (1960).

\begin{theorem}[completeness of resolution]\label{thm:rescomplete}
If a finite set $S$ of clauses is unsatisfiable, it has a refutation.
\end{theorem}
\begin{proof}
Induction on the number $k$ of variables occurring in $S$. If $k=0$, the only clause without variables is $\Box$, and the empty set
of clauses is satisfied by every assignment; so $S$, being unsatisfiable, contains $\Box$, a one-line refutation.

Let $k\ge1$ and choose a variable $x$ of $S$. Split $S$ into $S_0$, the clauses without $x$ or $\neg x$, $P$, the clauses containing $x$,
and $N$, those containing $\neg x$; a tautological clause containing both may be dropped, since it is satisfied by every
assignment and does not change satisfiability. Let $S'$ consist of the clauses of $S_0$ and of all resolvents on $x$ of a clause of $P$ and a clause of $N$. The
variable $x$ does not occur in $S'$.

We claim that $S'$ is unsatisfiable. Suppose $\alpha$ satisfies $S'$; we extend it to $x$ so as to satisfy $S$, a contradiction. Write
the clauses of $P$ as $C_i\cup\{x\}$ and those of $N$ as $D_j\cup\{\neg x\}$. If every $C_i$ is true under $\alpha$, set $x=0$: the clauses of $P$ hold
through $C_i$ and those of $N$ through $\neg x$. Otherwise some $C_i$ is false under $\alpha$; since every resolvent $C_i\cup D_j$ is true,
every $D_j$ is true, and setting $x=1$ satisfies $P$ through $x$ and $N$ through $D_j$. Either way $S_0$ holds, as $S_0\subseteq S'$.

So $S'$ is unsatisfiable with fewer variables, and by induction it has a refutation. Every clause of $S'$ is in $S$ or is a
resolvent of two clauses of $S$; prefixing those derivations gives a refutation of $S$.
\end{proof}

By Proposition~\ref{prop:refute}, resolution therefore proves every consequence: $S\models C$ if and only if $S\cup\{\{\bar\ell\}:\ell\in C\}$
has a refutation. But completeness says nothing about length. The \emph{pigeonhole formula} $\mathrm{PHP}_n$ has variables
$p_{ij}$ (pigeon $i\le n+1$ sits in hole $j\le n$), the clauses $p_{i1}\lor\dots\lor p_{in}$ (each pigeon sits somewhere) and $\neg p_{ij}\lor\neg p_{i'j}$ for
$i<i'$ (no two pigeons share a hole). It is unsatisfiable, and every human sees why at once; yet Haken (1985) proved
that every resolution refutation of $\mathrm{PHP}_n$ has at least $c^n$ clauses for some constant $c>1$. Finding short proofs, and proving
that none exist, is the subject of proof complexity, met again in Tome~C3.

\subsection*{DPLL}

Modern SAT solvers do not saturate a formula with resolvents; they search for a model and learn from failure. Their
skeleton is the \emph{DPLL} procedure (Davis, Putnam, Logemann and Loveland, 1962). For a set of clauses $S$ and a literal $\ell$,
let $S|_\ell$ be the result of making $\ell$ true: delete the clauses containing $\ell$ and remove $\bar\ell$ from the others.

\begin{center}
\begin{minipage}{0.9\linewidth}
$\mathrm{DPLL}(S)$:
\begin{enumerate}
\item If $S$ is empty, return \emph{satisfiable}. If $\Box\in S$, return \emph{unsatisfiable}.
\item \emph{Unit propagation}: if $S$ contains a one-literal clause $\{\ell\}$, return $\mathrm{DPLL}(S|_\ell)$.
\item \emph{Pure literal}: if a literal $\ell$ occurs in $S$ and $\bar\ell$ does not, return $\mathrm{DPLL}(S|_\ell)$.
\item \emph{Branch}: choose a literal $\ell$ of $S$; if $\mathrm{DPLL}(S|_\ell)$ is satisfiable return it, else return $\mathrm{DPLL}(S|_{\bar\ell})$.
\end{enumerate}
\end{minipage}
\end{center}

\begin{proposition}\label{prop:dpll}
DPLL terminates and answers correctly. Moreover the literals chosen along a successful branch, with the other variables
set arbitrarily, form a model of $S$.
\end{proposition}
\begin{proof}
Each recursive call is on a set of clauses with one variable fewer, so the recursion depth is at most the number of
variables. Correctness rests on three facts, for an assignment $\alpha$ with $\alpha(\ell)=1$: $\alpha$ satisfies $S$ if and only if it
satisfies $S|_\ell$ (the deleted clauses hold through $\ell$, and $\bar\ell$ is false anyway), and $S|_\ell$ does not involve the variable of $\ell$.
So $S$ is satisfiable if and only if $S|_\ell$ or $S|_{\bar\ell}$ is, which justifies step~4. In step~2 every model of $S$ makes $\ell$ true, so $S$ is
satisfiable if and only if $S|_\ell$ is. In step~3, if $\alpha$ is a model of $S$, the assignment $\alpha'$ that agrees with $\alpha$ except that it makes $\ell$
true is still a model: the clauses containing $\ell$ hold, and the others do not contain $\bar\ell$, so their truth is unchanged.
Hence $S$ is satisfiable if and only if $S|_\ell$ is. Step~1 is correct by definition. The last statement follows by
induction from the equivalence ``$\alpha$ satisfies $S$ iff it satisfies $S|_\ell$'' when $\alpha(\ell)=1$.
\end{proof}

A run of DPLL on an unsatisfiable formula is, read backwards, a tree-like resolution refutation (Exercise~\ref{ex:dpllres}), so
Haken's theorem bounds it too. Current solvers add \emph{clause learning}: when a branch fails, they derive by resolution a
new clause that explains the failure and keep it, which lets them solve industrial instances with millions of clauses.
The companion notebook builds a DPLL solver with unit propagation and watched literals and uses it throughout this book.

\section{Compactness}

Here is a puzzle about infinity. You are given infinitely many requirements, each about finitely many switches, and
someone has checked that any finite batch of them can be met at once. Can they all be met at once? It sounds like it
should need a miracle; it needs only patience. An infinite set of formulas, each mentioning finitely many variables, is satisfiable as soon as each of its finite
parts is. This is a theorem about infinity, and the proof is a careful choice made one variable at a time.

\begin{theorem}[compactness]\label{thm:compactness}
Let $\Gamma$ be a set of formulas in the countably many variables $x_1,x_2,\dots$. If every finite subset of $\Gamma$ is satisfiable, then
$\Gamma$ is satisfiable.
\end{theorem}
\begin{proof}
Call an assignment $\beta$ of $x_1,\dots,x_n$ \emph{good} if every finite subset of $\Gamma$ has a model extending $\beta$. The empty assignment
($n=0$) is good by hypothesis. If $\beta$ is good, one of its two extensions $\beta0$ and $\beta1$ (setting $x_{n+1}$ to $0$ or to $1$) is good. Otherwise there
are finite sets $F_0,F_1\subseteq\Gamma$ such that no model of $F_0$ extends $\beta0$ and no model of $F_1$ extends $\beta1$. But $F_0\cup F_1$ is finite, so it has a
model extending $\beta$, and that model extends $\beta0$ or $\beta1$ and satisfies both $F_0$ and $F_1$: a contradiction.

So we can choose $\varepsilon_1,\varepsilon_2,\ldots\in\{0,1\}$ one after the other such that each assignment $x_1\mapsto\varepsilon_1,\dots,x_n\mapsto\varepsilon_n$ is good. Let
$\alpha(x_i)=\varepsilon_i$ for all $i$. A formula $\varphi\in\Gamma$ only involves variables among $x_1,\dots,x_n$ for some $n$. The restriction of $\alpha$ to $x_1,\dots,x_n$
is good, so $\{\varphi\}$ has a model that agrees with $\alpha$ on $x_1,\dots,x_n$, and therefore $\alpha(\varphi)=1$.
\end{proof}

The same statement holds for any number of variables, but then the step-by-step choice must be replaced by Zorn's
lemma or by Tychonoff's theorem: the assignments form the compact space $\{0,1\}^X$ (Tome~6, Chapter~4), and each formula is
satisfied on a closed subset. The theorem gets its name from this.

\begin{corollary}[de Bruijn and Erd\H os, 1951]\label{cor:debruijnerdos}
A graph whose vertex set is countable can be coloured with $k$ colours, adjacent vertices getting different colours, if
and only if each of its finite subgraphs can.
\end{corollary}
\begin{proof}
One direction is clear. For the other, take a variable $x_{v,c}$ for each vertex $v$ and colour $c\in\{1,\dots,k\}$, meaning ``$v$ gets colour $c$''.
The set $\Gamma$ contains, for each vertex $v$, the clause $x_{v,1}\lor\dots\lor x_{v,k}$, and for each edge $vw$ and colour $c$, the clause $\neg x_{v,c}\lor\neg x_{w,c}$. A
finite subset of $\Gamma$ involves finitely many vertices; a colouring of the subgraph they span gives a model (setting $x_{v,c}$ true
when $v$ gets colour $c$). By compactness $\Gamma$ has a model $\alpha$; give $v$ the least colour $c$ with $\alpha(x_{v,c})=1$. Adjacent vertices then get
different colours.
\end{proof}

So the plane graph with all lattice points of $\Z^2$ joined at distance $1$ is $2$-colourable, a question about infinitely many
vertices answered by finite ones. Chapter~8 uses compactness to deduce the finite Ramsey theorem from the infinite one.

\section{Horn clauses and 2-SAT}\label{sec:twosat}

In general nobody knows how to decide satisfiability in polynomial time: that is the P versus NP question of Tome~C3.
But two natural fragments are easy, and the reasons are instructive.

\subsection*{Horn clauses}

Alfred Horn, an algebraist at UCLA, singled out these clauses in 1951 for a reason that had nothing to do with
computers (they are the sentences preserved by products of structures). Twenty years later they became the
programming language Prolog. A \emph{Horn clause} has at most one positive literal: $\neg a_1\lor\dots\lor\neg a_k\lor b$, that is, the rule $a_1\land\dots\land a_k\to b$ (a \emph{definite} clause; a
\emph{fact} if $k=0$), or $\neg a_1\lor\dots\lor\neg a_k$ (a \emph{goal}, $a_1\land\dots\land a_k\to\bot$). Rules of this shape are the language of logic
programming (Prolog), of database queries (Datalog) and of type inference.

\begin{theorem}\label{thm:horn}
Let $S$ be a finite set of Horn clauses. Start with all variables false and repeat: if a definite clause $a_1\land\dots\land a_k\to b$ has all
$a_i$ true and $b$ false, make $b$ true. When nothing changes, $S$ is satisfiable if and only if every goal clause has a false
variable, and then the final assignment $\mu$ is a model of $S$. Moreover $\mu$ is the least model: every model of $S$ makes true
every variable that $\mu$ makes true.
\end{theorem}
\begin{proof}
The loop changes at most one variable from false to true at each step, so it stops after at most as many steps as there
are variables. Let $\alpha$ be any model of $S$. By induction on the steps, every variable made true by the loop is true in $\alpha$: if
$a_1,\dots,a_k$ are true in $\alpha$, the clause $a_1\land\dots\land a_k\to b$ forces $b$ to be true in $\alpha$. So if some goal clause has all its variables
true in $\mu$, they are true in $\alpha$ too and $\alpha$ violates it: there is no model at all. Otherwise $\mu$ satisfies every goal
clause, and every definite clause too, because the loop stopped: a definite clause whose hypotheses are all true has
its conclusion true. So $\mu$ is a model, and the least one by the induction.
\end{proof}

With a counter per clause recording how many of its hypotheses are still false, the loop runs in time proportional
to the total length of $S$ (Dowling and Gallier, 1984); Tome~C2 analyses this kind of bookkeeping.

\subsection*{2-SAT}

Now let every clause have at most two literals; a one-literal clause $\{\ell\}$ is written $\ell\lor\ell$. The clause $a\lor b$ says $\bar a\to b$
and, equally, $\bar b\to a$. The \emph{implication graph} of $S$ has the $2n$ literals as vertices and, for each clause $a\lor b$, the two arrows
$\bar a\to b$ and $\bar b\to a$. It is \emph{skew-symmetric}: there is an arrow $u\to v$ if and only if there is an arrow $\bar v\to\bar u$, so there is a path
from $u$ to $v$ if and only if there is one from $\bar v$ to $\bar u$. Write $u\leadsto v$ for ``there is a path from $u$ to $v$''.

\begin{theorem}[Aspvall, Plass and Tarjan, 1979]\label{thm:twosat}
A set of $2$-clauses is satisfiable if and only if there is no variable $x$ with $x\leadsto\neg x$ and $\neg x\leadsto x$.
\end{theorem}
\begin{proof}
Every arrow $u\to v$ comes from a clause and is an implication valid in every model; hence so is every path. If $x\leadsto\neg x$
and $\neg x\leadsto x$, a model would make $x$ true (then $\neg x$ true, impossible) or false (then $\neg x$ true, and so $x$ true): there is none.

Conversely, assume the condition and induct on the number of variables. If $S$ is empty there is nothing to do. Otherwise
pick a variable $x$; by hypothesis $x\not\leadsto\neg x$ or $\neg x\not\leadsto x$, so there is a literal $\ell$ with $\ell\not\leadsto\bar\ell$. Make true every literal $m$
with $\ell\leadsto m$ (including $\ell$ itself). This is consistent: if $\ell\leadsto m$ and $\ell\leadsto\bar m$, skew symmetry turns the second path into
$m\leadsto\bar\ell$, and $\ell\leadsto m\leadsto\bar\ell$ is excluded. Every clause $a\lor b$ with a variable now assigned is satisfied: if $a$ was made true
there is nothing to say, and if $a$ was made false, then $\ell\leadsto\bar a$ and the arrow $\bar a\to b$ gives $\ell\leadsto b$, so $b$ was made true. The
clauses with no variable assigned form a set $S'$ whose implication graph is a subgraph of the old one, so $S'$ still
satisfies the condition and, by induction, has a model on the remaining variables. Together with the literals made
true, it is a model of $S$.
\end{proof}

\begin{figure}[ht]
\centering
\fig{ch01-implication}
\caption{The implication graph of the seven clauses of the opening question. The thick path $x_1\to x_3\to x_4\to\neg x_1$ and the
dashed path $\neg x_1\to x_2\to x_1$ put $x_1$ and $\neg x_1$ in the same strongly connected component: the seven clauses are
unsatisfiable. Without the clause $\neg x_2\lor x_1$, the two arrows $x_2\to x_1$ and $\neg x_1\to\neg x_2$ disappear.}
\label{fig:implication}
\end{figure}

The condition says that $x$ and $\neg x$ never lie in the same \emph{strongly connected component} of the implication graph, the
classes of the relation ``$u\leadsto v$ and $v\leadsto u$''. Tome~C2 computes these components in linear time (Tarjan's algorithm), so
2-SAT is solved in time proportional to the length of the formula, whereas 3-SAT is NP-complete. For the opening
question: the six clauses are satisfiable, for instance by $x_2=1$ and $x_1=x_3=x_4=0$ (the proof above finds such a model
starting from $\ell=x_2$), while the seven clauses are not (Figure~\ref{fig:implication}). Tome~3, Chapter~7 gives a
different, randomised algorithm for 2-SAT: a random walk on assignments.

\section{First-order logic}

Propositional logic cannot say ``every element has an inverse'' or ``between two rationals there is a third''. For
that we need variables that range over the elements of a structure, and quantifiers.

A \emph{signature} is a set of function symbols and relation symbols, each with an arity; constants are function symbols of
arity $0$. The signature of graphs has one binary relation $E$; that of ordered sets one binary relation $<$; that of
groups a binary function $\cdot$, a unary function ${}^{-1}$ and a constant $e$; that of arithmetic $0,1,+,\times,<$. \emph{Terms} are built from
variables $v_1,v_2,\dots$ by applying function symbols, and \emph{formulas} from atomic formulas $t_1=t_2$ and $R(t_1,\dots,t_k)$ by the
connectives and the quantifiers $\forall v$ and $\exists v$. A variable occurrence is \emph{bound} if it lies under a quantifier on that
variable, and \emph{free} otherwise; a \emph{sentence} has no free variable.

A \emph{structure} $\mathcal M$ for the signature is a non-empty set $M$ with a function $M^k\to M$ for each $k$-ary function symbol and a subset of
$M^k$ for each $k$-ary relation symbol. Given values in $M$ for the free variables, the truth of a formula in $\mathcal M$ is
defined by induction, the only new clauses being: $\forall v\,\varphi$ is true if $\varphi$ is true for every value of $v$ in $M$, and $\exists v\,\varphi$ if
it is true for some value (Tarski, 1933). We write $\mathcal M\models\varphi$ when a sentence $\varphi$ is true in $\mathcal M$, and $\Gamma\models\varphi$ when $\varphi$ is
true in every structure where all sentences of $\Gamma$ are.

For instance $\forall v_1\forall v_2\,(v_1<v_2\to\exists v_3\,(v_1<v_3\land v_3<v_2))$ is true in $(\Q,<)$ and false in $(\Z,<)$; and the sentence
$\forall v_1\exists v_2\,(v_1\cdot v_2=e)$ holds in every group.

\begin{proposition}[prenex form]\label{prop:prenex}
Every formula is equivalent (true for the same values of its free variables in every structure) to a formula
$Q_1w_1\cdots Q_kw_k\,\psi$, where the $Q_i$ are quantifiers and $\psi$ has no quantifier.
\end{proposition}
\begin{proof}
First rename bound variables so that each quantifier uses its own variable, distinct from the free variables and from
the other quantified ones; this does not change the meaning, as the name of a bound variable plays no role (an
induction on formulas). Remove $\to$ and $\leftrightarrow$ by their expressions in $\neg,\land,\lor$. Then move quantifiers outwards with the
equivalences $\neg\forall w\,\varphi\equiv\exists w\,\neg\varphi$, $\neg\exists w\,\varphi\equiv\forall w\,\neg\varphi$, and, when $w$ is not free in $\chi$,
$(\forall w\,\varphi)\land\chi\equiv\forall w\,(\varphi\land\chi)$, $(\exists w\,\varphi)\land\chi\equiv\exists w\,(\varphi\land\chi)$ and the same with $\lor$. Each of these is checked directly on the
definition of truth; the hypothesis that $w$ is not free in $\chi$ (secured by the renaming) is used, and so is the fact that
$M$ is non-empty, for instance in $(\forall w\,\varphi)\lor\chi\equiv\forall w\,(\varphi\lor\chi)$. Formally the proof is an induction on the formula, which
pulls the prefixes of the two sides of $\land$ or $\lor$ out one after the other.
\end{proof}

\begin{remark}
\emph{Gödel's completeness theorem} (1929) says that first-order consequence is captured by a finite proof system: if
$\Gamma\models\varphi$, there is a formal proof of $\varphi$ from $\Gamma$, which a computer can check. It implies the compactness theorem for
first-order logic: if every finite part of $\Gamma$ has a model, so does $\Gamma$. Both are proved in Tome~C3, Chapter~4, together with the
\emph{incompleteness} theorems, which say that for arithmetic the true sentences cannot all be captured by any such
system, and with Church and Turing's theorem that no algorithm decides whether a first-order sentence is valid. The
propositional theorems of this chapter are the shadows of these.
\end{remark}

\section{Exercises}

\begin{exercise}[$\star$]\label{ex:connectives}
Show that $\{\neg,\to\}$ is complete, and that $\{\land,\lor,\to,\leftrightarrow\}$ is not. Is $\{\to,\bot\}$ complete?
\end{exercise}

\begin{exercise}[$\star\star$]\label{ex:monotone}
Show that a Boolean function is monotone if and only if it is computed by a formula using only $\land$ and $\lor$ (or is a
constant). Give such a formula for the \emph{majority} of three bits.
\end{exercise}

\begin{exercise}[$\star\star$]\label{ex:shannoncount}
\emph{Most functions need long formulas.} Show that there are at most $(n+3)^{2s+2}$ formulas in $x_1,\dots,x_n$ with at most $s$
connectives among $\neg,\land,\lor$. Deduce that if $(2s+2)\log_2(n+3)<2^n$, some Boolean function of $n$ variables has no formula with at
most $s$ connectives. For $n=20$, how large is the bound?
\end{exercise}

\begin{exercise}[$\star$]\label{ex:lattice}
In a Boolean algebra, show that $x\land y$ is the greatest lower bound of $x$ and $y$ for the order of Lemma~\ref{lem:boolbasic}, and
$x\lor y$ the least upper bound.
\end{exercise}

\begin{exercise}[$\star\star$]\label{ex:php}
Write the clauses of $\mathrm{PHP}_2$ (three pigeons, two holes) and find a resolution refutation. How many clauses does your
refutation use?
\end{exercise}

\begin{exercise}[$\star$]\label{ex:threesat}
Show that a clause $\ell_1\lor\dots\lor\ell_k$ with $k\ge4$ can be replaced by the $k-2$ clauses
$\ell_1\lor\ell_2\lor y_1$, $\bar y_1\lor\ell_3\lor y_2$, \dots, $\bar y_{k-3}\lor\ell_{k-1}\lor\ell_k$, with new variables $y_i$, without changing satisfiability.
\end{exercise}

\begin{exercise}[$\star\star$]\label{ex:xorsat}
A \emph{parity constraint} says that $x_{i_1}\oplus\dots\oplus x_{i_k}=b$. Show how to decide in polynomial time whether a set of parity
constraints is satisfiable and to count its solutions. Why does this not contradict Proposition~\ref{prop:parity}?
\end{exercise}

\begin{exercise}[$\star\star$]\label{ex:hornmodels}
Show that the set of models of a Horn formula is closed under intersection: if $\alpha$ and $\beta$ are models, so is $\alpha\land\beta$ (the
assignment $x\mapsto\min(\alpha(x),\beta(x))$). Find a $2$-clause formula whose models are not closed under intersection.
\end{exercise}

\begin{exercise}[$\star\star\star$]\label{ex:hornconverse}
Conversely, show that if the set of models of a formula $\varphi$ is closed under intersection, then $\varphi$ is equivalent to a
conjunction of Horn clauses (McKinsey, 1943).
\end{exercise}

\begin{exercise}[$\star\star$]\label{ex:konig}
\emph{K\H onig's lemma.} A rooted tree in which every vertex has finitely many children and which has infinitely many
vertices has an infinite path from the root. Prove it directly, then deduce it from Theorem~\ref{thm:compactness}.
\end{exercise}

\begin{exercise}[$\star\star$]\label{ex:interpolation}
\emph{Craig interpolation.} Let $\varphi\to\psi$ be a tautology, and let $V$ be the set of variables common to $\varphi$ and $\psi$. Show
that there is a formula $\theta$ with variables in $V$ such that $\varphi\to\theta$ and $\theta\to\psi$ are tautologies. (Model checkers use
interpolants to over-approximate the reachable states of a program.)
\end{exercise}

\begin{exercise}[$\star\star$]\label{ex:dpllres}
Show that if DPLL without the pure literal rule answers \emph{unsatisfiable} on $S$, then $S$ has a resolution refutation with
at most one clause per node of the recursion tree, plus the clauses of $S$.
\end{exercise}

\begin{exercise}[$\star$]\label{ex:fosentences}
Write first-order sentences saying, in the signature of graphs: (a) the graph has no triangle; (b) every vertex has
exactly two neighbours; (c) the graph has a vertex adjacent to all others. In the signature $\{<\}$, write a sentence true
in $(\R,<)$ and false in $(\Z,<)$, and one true in $(\N,<)$ and false in $(\Z,<)$.
\end{exercise}

\section{Hints}

\begin{hint}{ex:connectives}
Express $\lor$ with $\neg$ and $\to$. For the second set, look at the assignment that makes every variable true. For the
third, express $\neg$.
\end{hint}
\begin{hint}{ex:monotone}
Take the disjunction, over the minimal points $a$ where $f(a)=1$, of the conjunction of the variables $x_i$ with $a_i=1$.
\end{hint}
\begin{hint}{ex:shannoncount}
Write formulas in Polish notation (connective first), which needs no parentheses. How many variable occurrences can a
formula with $s$ connectives have?
\end{hint}
\begin{hint}{ex:lattice}
For the lower bound: if $z\le x$ and $z\le y$, compute $z\land(x\land y)$. For the upper bound, use Lemma~\ref{lem:boolbasic}(d).
\end{hint}
\begin{hint}{ex:php}
First derive the one-literal clause $\{\neg p_{21}\}$ (pigeon $2$ is not in hole $1$) and the clause $\{p_{21}\}$, each by
following where pigeons $1$ and $3$ could go.
\end{hint}
\begin{hint}{ex:threesat}
If $\ell_j$ is true, make $y_m$ true for the clauses before the one containing $\ell_j$, false after it.
\end{hint}
\begin{hint}{ex:xorsat}
A parity constraint is a linear equation over $\mathbb F_2$.
\end{hint}
\begin{hint}{ex:hornmodels}
If $\alpha\land\beta$ makes the hypotheses of a rule true, so do $\alpha$ and $\beta$.
\end{hint}
\begin{hint}{ex:hornconverse}
Let $H$ be the conjunction of all Horn clauses implied by $\varphi$. Given an assignment $\beta$ that is not a model of $\varphi$,
intersect all the models of $\varphi$ that lie above $\beta$.
\end{hint}
\begin{hint}{ex:konig}
Directly: walk down, always to a child with infinitely many descendants. With compactness: one variable per vertex,
saying ``the path goes through $v$''.
\end{hint}
\begin{hint}{ex:interpolation}
Induction on the number of variables of $\varphi$ that are not in $V$. To remove such a variable $x$, replace $\varphi$ by
$\varphi[x:=\top]\lor\varphi[x:=\bot]$.
\end{hint}
\begin{hint}{ex:dpllres}
Show by induction, from the leaves up, that at a node where the literals of a set $\rho$ have been made true, there is a
derived clause all of whose literals are made false by $\rho$.
\end{hint}
\begin{hint}{ex:fosentences}
``Exactly two'' means: there are two distinct neighbours, and every neighbour is one of them.
\end{hint}

\section{Solutions}

\begin{solution}{ex:connectives}
\idea To show completeness, rebuild a set already known to be complete. To show incompleteness, find a property that every
formula over the set has and some function lacks, as in Proposition~\ref{prop:monotone}.
\writeup $\varphi\lor\psi\equiv\neg\varphi\to\psi$, and $\{\neg,\lor\}$ is complete since $\varphi\land\psi\equiv\neg(\neg\varphi\lor\neg\psi)$; so $\{\neg,\to\}$ is complete. Let $\mathbf 1$ be the assignment
making every variable true. By structural induction, every formula built from variables with $\land,\lor,\to,\leftrightarrow$ is true at $\mathbf 1$,
because each of these connectives gives $1$ on two inputs equal to $1$. So $\neg x_1$, false at $\mathbf 1$, is not expressible. Finally
$\neg\varphi\equiv\varphi\to\bot$, so $\{\to,\bot\}$ can express $\neg$ and $\to$, and is complete.
\end{solution}

\begin{solution}{ex:monotone}
\idea A monotone function is determined by its minimal ``true'' points: it is $1$ exactly above one of them, and ``above $a$'' is
a conjunction of variables.
\writeup Proposition~\ref{prop:monotone} gives one direction. Conversely let $f$ be monotone and not constant, and let $A$ be the set of
minimal points $a$ with $f(a)=1$ (for the coordinatewise order). Since $f(0,\dots,0)=0$ (otherwise $f\equiv1$ by monotonicity), every
$a\in A$ has a coordinate $1$, and $t_a=\bigwedge_{a_i=1}x_i$ is a formula without negation, true at $b$ if and only if $a\le b$. Let
$\varphi=\bigvee_{a\in A}t_a$. If $f(b)=1$, there is a minimal $a\le b$ with $f(a)=1$ (go down while possible), and $t_a$ is true at $b$. If $\varphi$ is
true at $b$, then $a\le b$ for some $a\in A$, and $f(b)\ge f(a)=1$. So $\varphi$ computes $f$. For majority, the minimal true points are
$110,101,011$, giving $(x_1\land x_2)\lor(x_1\land x_3)\lor(x_2\land x_3)$.
\end{solution}

\begin{solution}{ex:shannoncount}
\idea Count formulas by writing them as strings, and compare with the $2^{2^n}$ functions: if there are fewer short formulas
than functions, some function has no short formula. This is Shannon's counting argument (1949), used again in Tome~C3
for circuits.
\writeup In Polish notation a formula is written as its connective followed by its arguments, for example $\land\,x_1\neg x_2$; it is a
string over the $n+3$ symbols $x_1,\dots,x_n,\neg,\land,\lor$, and different parse trees give different strings (reading from the
left, the first symbol tells how many arguments follow, and an induction shows that a formula has a unique reading).
A formula with $c\le s$ connectives, of which $b$ are binary, has $b+1\le s+1$ leaves, so its length is at most $2s+1$. The number of
strings of length at most $2s+1$ is $\sum_{L=0}^{2s+1}(n+3)^L\le(n+3)^{2s+2}$ (a geometric sum with ratio at least $2$). There are $2^{2^n}$
functions, and $(n+3)^{2s+2}<2^{2^n}$ if and only if $(2s+2)\log_2(n+3)<2^n$. For $n=20$ this holds for every $s\le115\,900$: some
function of $20$ variables needs more than $115\,900$ connectives, while its truth table has about a million bits. In
fact almost all functions do, since the short formulas cover a vanishing fraction of the functions.
\end{solution}

\begin{solution}{ex:lattice}
\idea Translate ``$\le$'' into equations with $\land$ and use associativity.
\writeup $(x\land y)\land x=x\land y$, so $x\land y\le x$, and likewise $x\land y\le y$. If $z\le x$ and $z\le y$, then $z\land(x\land y)=(z\land x)\land y=z\land y=z$, so
$z\le x\land y$. By Lemma~\ref{lem:boolbasic}(d), $u\le v$ means $u\lor v=v$. Since $x\lor(x\lor y)=x\lor y$, $x\le x\lor y$, and likewise $y\le x\lor y$. If $x\le z$
and $y\le z$, then $(x\lor y)\lor z=x\lor(y\lor z)=x\lor z=z$, so $x\lor y\le z$.
\end{solution}

\begin{solution}{ex:php}
\idea Reason as a human would, ``if pigeon $2$ were in hole $1$, then pigeons $1$ and $3$ would both have to be in hole $2$'', and turn each
step into a resolution.
\writeup The clauses are $\{p_{i1},p_{i2}\}$ for $i=1,2,3$ and $\{\neg p_{ij},\neg p_{kj}\}$ for $j=1,2$ and $i<k$: nine clauses. A refutation
with $11$ resolutions (the variable resolved on is in brackets):
\begin{enumerate}
\item $\{p_{11},p_{12}\}$ and $\{\neg p_{12},\neg p_{32}\}$ give $\{p_{11},\neg p_{32}\}$ $[p_{12}]$;
\item $\{p_{31},p_{32}\}$ and $\{\neg p_{21},\neg p_{31}\}$ give $\{\neg p_{21},p_{32}\}$ $[p_{31}]$;
\item from 1 and 2: $\{p_{11},\neg p_{21}\}$ $[p_{32}]$; with $\{\neg p_{11},\neg p_{21}\}$: $\{\neg p_{21}\}$ $[p_{11}]$;
\item $\{p_{21},p_{22}\}$ and $\{\neg p_{12},\neg p_{22}\}$ give $\{\neg p_{12},p_{21}\}$ $[p_{22}]$; with $\{p_{11},p_{12}\}$: $\{p_{11},p_{21}\}$ $[p_{12}]$;
\item $\{p_{21},p_{22}\}$ and $\{\neg p_{22},\neg p_{32}\}$ give $\{p_{21},\neg p_{32}\}$ $[p_{22}]$; $\{p_{31},p_{32}\}$ and $\{\neg p_{11},\neg p_{31}\}$ give $\{\neg p_{11},p_{32}\}$ $[p_{31}]$;
these two give $\{\neg p_{11},p_{21}\}$ $[p_{32}]$;
\item from the last clauses of 4 and 5: $\{p_{21}\}$ $[p_{11}]$; and with $\{\neg p_{21}\}$ from 3: $\Box$ $[p_{21}]$.
\end{enumerate}
The notebook finds this refutation by breadth-first search and checks that the saturation of the nine clauses
contains $73$ non-tautological clauses.
\end{solution}

\begin{solution}{ex:threesat}
\idea The new variables form a chain of ``the clause is not yet satisfied'' flags passed from one short clause to the next.
\writeup Number the new clauses $C_1=\ell_1\lor\ell_2\lor y_1$, $C_m=\bar y_{m-1}\lor\ell_{m+1}\lor y_m$ for $2\le m\le k-3$, and $C_{k-2}=\bar y_{k-3}\lor\ell_{k-1}\lor\ell_k$. If all $\ell_i$
are false in an assignment, then $C_1$ forces $y_1$, then $C_2$ forces $y_2$, and so on up to $y_{k-3}$, and $C_{k-2}$ is false. So every model of the
new clauses satisfies the old clause. Conversely, if some $\ell_j$ is true, let $C_c$ be the new clause containing it and set
$y_m=1$ for $m<c$ and $y_m=0$ for $m\ge c$. Then $C_m$ holds through $y_m$ for $m<c$, $C_c$ through $\ell_j$, and $C_m$ through $\bar y_{m-1}$ for $m>c$.
Doing this for every long clause, with its own new variables, turns any CNF into a $3$-CNF with the same satisfiability.
\end{solution}

\begin{solution}{ex:xorsat}
\idea Over $\mathbb F_2$, $\oplus$ is addition, so the constraints form a linear system.
\writeup Each constraint is the equation $x_{i_1}+\dots+x_{i_k}=b$ in $\mathbb F_2$. Gaussian elimination on the $m\times n$ system, in $O(mn\min(m,n))$
operations, decides consistency; if the system is consistent with rank $r$, its solutions form an affine space of
dimension $n-r$, with $2^{n-r}$ points. There is no contradiction with Proposition~\ref{prop:parity}: that proposition bounds the
size of a CNF \emph{equivalent} to parity, without new variables. With new variables $y_1,\dots,y_{k-1}$ and the constraints
$y_1=x_1\oplus x_2$, $y_{i}=y_{i-1}\oplus x_{i+1}$, each a CNF of four clauses, parity gets an equisatisfiable CNF of linear size, as in
Theorem~\ref{thm:tseitin}. But a SAT solver that only resolves does not ``see'' linear algebra: Tseitin's parity formulas on
expander graphs need exponentially long resolution refutations (Urquhart, 1987), which is why solvers add Gaussian
elimination for XOR constraints.
\end{solution}

\begin{solution}{ex:hornmodels}
\idea A Horn clause says ``if all these, then that''; an intersection can only make fewer hypotheses true.
\writeup Let $\alpha,\beta$ be models and $\gamma=\alpha\land\beta$. For a rule $a_1\land\dots\land a_k\to b$: if $\gamma$ makes all $a_i$ true, then so do $\alpha$ and $\beta$, so both make $b$
true, and so does $\gamma$. For a goal $\neg a_1\lor\dots\lor\neg a_k$: if $\gamma$ made all $a_i$ true, so would $\alpha$, which is a model. The clause
$x\lor y$ has the models $(1,0)$ and $(0,1)$, whose intersection $(0,0)$ is not a model.
\end{solution}

\begin{solution}{ex:hornconverse}
\idea To show that $\varphi$ is equivalent to the conjunction $H$ of its Horn consequences, we must separate each non-model $\beta$ from
the models by a Horn clause false at $\beta$. Closure under intersection gives a least model above $\beta$, which tells us which
conclusion to draw.
\writeup Let $x_1,\dots,x_n$ be the variables of $\varphi$; all assignments below are of these variables. Let $H$ be the conjunction of the
(finitely many) Horn clauses in $x_1,\dots,x_n$ that are consequences of $\varphi$. Then $\varphi\models H$, and it remains to show that every
model of $H$ is a model of $\varphi$, that is, that every $\beta$ which is not a model of $\varphi$ falsifies some clause of $H$. Let
$T=\{i:\beta(x_i)=1\}$, and let $U$ be the set of models $\alpha$ of $\varphi$ with $\alpha\ge\beta$ (that is, true on $T$).

If $U$ is empty, the goal clause $\bigvee_{i\in T}\neg x_i$ is a consequence of $\varphi$ (every model makes some $x_i$, $i\in T$, false), and $\beta$ falsifies it.

Otherwise, let $\gamma$ be the intersection of the members of $U$; by hypothesis it is a model of $\varphi$, and $\gamma\ge\beta$. As $\beta$ is not a model, $\gamma\ne\beta$,
so there is $j$ with $\gamma(x_j)=1$ and $\beta(x_j)=0$. The Horn clause $\bigwedge_{i\in T}x_i\to x_j$ is a consequence of $\varphi$: a model that makes all $x_i$,
$i\in T$, true belongs to $U$, so it lies above $\gamma$ and makes $x_j$ true. And $\beta$ falsifies it. So $\varphi\equiv H$.
\end{solution}

\begin{solution}{ex:konig}
\idea An infinite set split into finitely many parts has an infinite part: the pigeonhole principle, applied at each step
down the tree.
\writeup \emph{Directly.} Call a vertex \emph{rich} if it has infinitely many descendants. The root is rich. A rich vertex has finitely
many children, and its descendants are the children and their descendants, so at least one child is rich. Starting
from the root and always moving to a rich child gives an infinite path.

\emph{By compactness.} The tree has finitely many vertices at each depth (induction on the depth), so countably many in all.
Take a variable $x_v$ per vertex. Let $\Gamma$ contain, for each depth $n$, a formula saying that exactly one vertex of depth $n$ is
chosen (a finite formula, as depth $n$ is finite), and for each vertex $v\ne$ root the formula $x_v\to x_{\text{parent}(v)}$. A finite subset of
$\Gamma$ mentions depths at most $N$. Depth $N$ is not empty, otherwise all vertices would have depth less than $N$ and the tree
would be finite; choose $w$ at depth $N$ and make true exactly the $x_v$ for $v$ on the path from the root to $w$: this satisfies
the finite subset. By compactness, $\Gamma$ has a model. It chooses one vertex $v_n$ at each depth $n$, and the parent of $v_{n+1}$ is
chosen, so it is $v_n$: the $v_n$ form an infinite path.
\end{solution}

\begin{solution}{ex:interpolation}
\idea A variable that $\psi$ does not see can be ``projected out'' of $\varphi$: $\varphi[x:=\top]\lor\varphi[x:=\bot]$ says ``some value of $x$ makes $\varphi$ true''.
It is weaker than $\varphi$ and still strong enough for $\psi$, which does not care about $x$.
\writeup Induction on the number $k$ of variables of $\varphi$ outside $V$. If $k=0$, take $\theta=\varphi$. Otherwise let $x$ be such a variable and
$\varphi'=\varphi[x:=\top]\lor\varphi[x:=\bot]$, which has $k-1$ variables outside $V$. Then $\varphi\to\varphi'$ is a tautology: under any assignment $\alpha$, $\varphi$ has the
value of one of the two disjuncts, the one with the constant $\alpha(x)$. And $\varphi'\to\psi$ is a tautology: if $\alpha$ makes $\varphi'$ true, then
$\varphi$ is true at the assignment $\alpha'$ equal to $\alpha$ except perhaps at $x$; then $\psi$ is true at $\alpha'$, and since $x$ does not occur in
$\psi$, at $\alpha$. The variables common to $\varphi'$ and $\psi$ are in $V$, so by induction there is $\theta$ over $V$ with $\varphi'\to\theta$ and $\theta\to\psi$
tautologies; then $\varphi\to\theta$ as well. (If $V$ is empty, $\theta$ is a constant: $\varphi$ is unsatisfiable or $\psi$ is a tautology.) The
interpolant can be exponentially long; computing small interpolants from resolution proofs is what model checkers do.
\end{solution}

\begin{solution}{ex:dpllres}
\idea A failed branch on $\ell$ and a failed branch on $\bar\ell$ each produce a reason; resolving the two reasons on $\ell$ produces a
reason for the parent's failure.
\writeup For a node of the recursion tree, let $\rho$ be the set of literals made true on the way to it, so the node runs DPLL on $S|_\rho$.
We show by induction from the leaves up that if the node answers \emph{unsatisfiable}, some clause $C$ derivable from $S$, using one
new clause per node below it, has all its literals false under $\rho$ (we say $C$ is \emph{falsified} by $\rho$). At a leaf, $\Box\in S|_\rho$, which
means that a clause $C$ of $S$ is falsified by $\rho$. At a branching node on $\ell$, the children give clauses $C_1$ falsified by $\rho\cup\{\ell\}$ and $C_2$
falsified by $\rho\cup\{\bar\ell\}$. If $\bar\ell\notin C_1$, then $C_1$ is falsified by $\rho$ already; likewise if $\ell\notin C_2$. Otherwise the resolvent of $C_1$ and
$C_2$ on $\ell$ is falsified by $\rho$. At a unit propagation node, the unit clause $\{\ell\}$ of $S|_\rho$ comes from a clause $D\in S$ whose literals other
than $\ell$ are falsified by $\rho$; the child gives $C_1$ falsified by $\rho\cup\{\ell\}$, and either $C_1$ or its resolvent with $D$ on $\ell$ works. At the
root $\rho=\emptyset$, and the only clause falsified by the empty assignment is $\Box$. The refutation is tree-like and has at most one
derived clause per node.
\end{solution}

\begin{solution}{ex:fosentences}
\idea Translate word by word, using $=$ to express ``distinct'' and ``one of'', and remember that a graph here has a symmetric
irreflexive relation $E$, so the vertices of a triangle are automatically distinct.
\writeup (a) $\neg\exists a\,\exists b\,\exists c\,(E(a,b)\land E(b,c)\land E(a,c))$. (b) $\forall a\,\exists b\,\exists c\,\bigl(b\ne c\land E(a,b)\land E(a,c)\land\forall d\,(E(a,d)\to(d=b\lor d=c))\bigr)$.
(c) $\exists a\,\forall b\,(b\ne a\to E(a,b))$. Density, $\forall a\,\forall b\,(a<b\to\exists c\,(a<c\land c<b))$, is true in $\R$ and false in $\Z$ (nothing
lies strictly between $0$ and $1$). The existence of a least element, $\exists a\,\forall b\,\neg(b<a)$, is true in $\N$ and false in $\Z$.
\end{solution}

\section*{Sources and further reading}

\begin{itemize}
\item \emph{The stories.} The epigraph is Leibniz in B.~Russell's translation, \emph{A Critical Exposition of the Philosophy of Leibniz}, Cambridge University Press, 1900, \S\,85, shortened; the Latin
is in C.~I.~Gerhardt (ed.), \emph{Die philosophischen Schriften von G.~W.~Leibniz}, vol.~VII, 1890. Boole's white sheep are in G.~Boole, \emph{An Investigation of the Laws of Thought}, Walton and Maberly, 1854,
Chapter~II; his life is told in D.~MacHale, \emph{The Life and Work of George Boole}, Cork University Press, 2014.
\item H.~B.~Enderton, \emph{A Mathematical Introduction to Logic}, 2nd ed., Academic Press, 2001. Chapter~1 for propositional
logic and compactness, Chapter~2 for first-order structures and prenex forms.
\item D.~E.~Knuth, \emph{The Art of Computer Programming}, Volume~4B, \emph{Combinatorial Algorithms, Part~2}, Addison-Wesley,
2022. Section~7.2.2.2, \emph{Satisfiability}: DPLL, watched literals, clause learning, and a wealth of encodings.
\item A.~Biere, M.~Heule, H.~van Maaren and T.~Walsh (eds.), \emph{Handbook of Satisfiability}, 2nd ed., IOS Press, 2021.
\item M.~Davis and H.~Putnam, \emph{A computing procedure for quantification theory}, Journal of the ACM 7 (1960),
201--215; M.~Davis, G.~Logemann and D.~Loveland, \emph{A machine program for theorem-proving}, Communications of the
ACM 5 (1962), 394--397.
\item J.~A.~Robinson, \emph{A machine-oriented logic based on the resolution principle}, Journal of the ACM 12 (1965),
23--41.
\item G.~S.~Tseitin, \emph{On the complexity of derivation in propositional calculus}, in \emph{Studies in Constructive
Mathematics and Mathematical Logic}, Part~II, 1968, 115--125.
\item A.~Haken, \emph{The intractability of resolution}, Theoretical Computer Science 39 (1985), 297--308.
\item B.~Aspvall, M.~F.~Plass and R.~E.~Tarjan, \emph{A linear-time algorithm for testing the truth of certain quantified
Boolean formulas}, Information Processing Letters 8 (1979), 121--123. The 2-SAT theorem.
\item W.~F.~Dowling and J.~H.~Gallier, \emph{Linear-time algorithms for testing the satisfiability of propositional Horn
formulae}, Journal of Logic Programming 1 (1984), 267--284.
\item M.~H.~Stone, \emph{The theory of representations for Boolean algebras}, Transactions of the AMS 40 (1936), 37--111.
\item N.~G.~de Bruijn and P.~Erd\H os, \emph{A colour problem for infinite graphs and a problem in the theory of relations},
Indagationes Mathematicae 13 (1951), 371--373.
\item C.~E.~Shannon, \emph{The synthesis of two-terminal switching circuits}, Bell System Technical Journal 28 (1949),
59--98. The counting argument of Exercise~\ref{ex:shannoncount}.
\end{itemize}
